
\subsection{Noether's 定理}
\begin{tcolorbox}[colback=white!5,colframe=black!45]
          \centering
      \textbf{连续的对称性} 
      \(\longrightarrow \)\textbf{守恒律}
\end{tcolorbox}
\noindent
\\
场论中的Noether's theorem 推导有一些细节(\text{\color{blue}比如测度的影响,场的影响,坐标的影响}),推导过程如下：
\medskip
        $$
        x^\mu \rightarrow x'^\mu = x^\mu + \delta x^\mu(x)
        $$
        $$
        \Psi_\ell(x) \rightarrow \Psi_\ell (x) = \Psi_\ell(x) + \delta \Psi_\ell(x)
        $$
\noindent
场和坐标的变换带来如下两种影响：
\begin{itemize}
  \item \textbf{两种变分的变化}\\
            \begin{enumerate}
              \item \textbf{坐标固定的边分}:
                \(x^\mu\) 是固定的,仅是场的变化：
                \[\delta\Psi_\ell(x)=\Psi_\ell (x) - \Psi_\ell(x)\]
              \item \textbf{全变分(这里包含了场和坐标的共同影响)}这是指在无穷小坐标变换 
                \(x^\mu \rightarrow x'^\mu = x^\mu + \delta x^\mu(x)\)下，坐标变化与场的变化共同作用所导致的场的总变化:
                    \begin{align*}
                \delta_{total} \Psi_\ell(x)&= \Psi_\ell (x') - \Psi_\ell(x) \\
                &= \Psi_\ell (x + \delta x) - \Psi_\ell(x) \\
                &= \Psi_\ell (x) + \left(\partial_\nu \Psi_\ell (x)\right) \delta x^\nu - \Psi_\ell(x) \\
                &= \left(\Psi_\ell(x) + \delta \Psi_\ell(x)\right) + \partial_\nu \left(\Psi_\ell(x) + \delta \Psi_\ell(x)\right) \delta x^\nu - \Psi_\ell(x) \\
                &= \delta \Psi_\ell(x) + \left(\partial_\nu \Psi_\ell(x) \right)\delta  x^\nu+ \mathscr{O}(\delta^2)
                \end{align*}
                区分这两种形式对于研究对称变换和推导诺特流至关重要,因为它同时考虑了场变量和时空坐标的变化。
            \end{enumerate}
  \item \textbf{积分测度的变化}
      \[
      d^4 x' = \det\left(\frac{\partial x'^\nu}{\partial x^\mu}\right) d^4 x = (1 + \partial_\mu \delta x^\mu) d^4 x 
      \]
\end{itemize}

\(\implies\)\textbf{他们共同导致了作用量的变化：}\footnote{推导中，我加入了一个全导数项，它并不影响运动方程的变化，这实际上是规范项。在其他地方好像是诺特第二定理，我并不懂数学，这里暂且放在一起} 
\begin{equation}
    \begin{aligned}\label{3.2}
          \Delta I &=I[\Psi_\ell (x'),x']-I[\Psi_\ell(x),x]+\underbrace{\int d^4x\partial_\mu F^\mu}_\text{这实际上是规范项}\\
          &= \int d^4 x' \, \mathscr{L}'[\Psi_\ell (x'), \partial_\mu \Psi_\ell (x')] - \int d^4 x \, \mathscr{L}[\Psi_\ell(x), \partial_\mu \Psi_\ell(x)]+\int d^4x\partial_\mu F^\mu \\
          &= \int d^4 x (1 + \partial_\mu \delta x^\mu) \underbrace{\mathscr{L}'[\Psi_\ell (x'), \partial_\mu \Psi_\ell (x')]}_{\text{所有的变化都放在差值项里}\Delta\mathscr{L}} - \int d^4 x \, \mathscr{L}[\Psi_\ell(x), \partial_\mu \Psi_\ell(x)] +\int d^4x\partial_\mu F^\mu \\
          &= \int d^4 x (1 + \partial_\mu \delta x^\mu) \left\{\mathscr{L}[\Psi_\ell(x), \partial_\mu \Psi_\ell(x)]+\Delta\mathscr{L}\right\} - \int d^4 x \, \mathscr{L}[\Psi_\ell(x), \partial_\mu \Psi_\ell(x)]  +\int d^4x\partial_\mu F^\mu\\          
          &= \int d^4 x \left\{ \mathscr{L}[\Psi_\ell, \partial_\mu \Psi_\ell]+\Delta\mathscr{L} + (\partial_\mu \delta x^\mu) \mathscr{L}[\Psi_\ell , \partial_\mu \Psi_\ell ]  +\mathscr{O}(\varphi^{(a)})- \mathscr{L}[\Psi_\ell, \partial_\mu \Psi_\ell]\right\} +\int d^4x\partial_\mu F^\mu\\
          &= \int d^4 x \left\{ \underbrace{\Delta\mathscr{L}}_{\mathbf{1}}
          + \underbrace{(\partial_\mu \delta x^\mu )\mathscr{L}[\Psi_\ell (x), \partial_\mu \Psi_\ell (x)]}_{\mathbf{2}} +\partial_\mu F^\mu\right\} \\
    \end{aligned}
\end{equation}
上式(\refeq{3.2}),细节的来处理一些东西：
\begin{itemize}
  \item \textbf{第1项\(\Delta\mathscr{L}\):}
值得注意的是，由于我们已经固定（或统一）了积分测度，因此这里应当采用在固定坐标下进行的变分。
    \begin{align*}
      \Delta \mathscr{L} &= \mathscr{L}[\Psi_\ell , \partial_\mu \Psi_\ell ] - \mathscr{L}[\Psi_\ell, \partial_\mu \Psi_\ell] \\
        &= \mathscr{L}[\Psi_\ell+\delta\Psi_\ell, \partial_\mu \Psi_\ell+\delta\Psi_\ell] - \mathscr{L}[\Psi_\ell, \partial_\mu \Psi_\ell] \\
        &=\frac{\partial\mathscr{L}}{\partial \Psi_\ell}\delta\Psi_\ell+\frac{\partial \mathscr{L}}{\partial (\partial_\mu \Psi_\ell)}\partial_\mu\delta\Psi_\ell\\
        &=\left[\frac{\partial\mathscr{L}}{\partial \Psi_\ell}-\frac{\partial \mathscr{L}}{\partial (\partial_\mu \Psi_\ell)} \right]\delta\Psi_\ell+\partial_\mu\left(\frac{\partial \mathscr{L}}{\partial (\partial_\mu \Psi_\ell)}\delta\Psi_\ell\right)
    \end{align*}
  \item \textbf{第2项 \(\mathscr{L}[\Psi_\ell (x), \partial_\mu \Psi_\ell (x)]\):}
    \begin{align*}
      \mathscr{L}[\Psi_\ell (x), \partial_\mu \Psi_\ell (x)] 
      &= \mathscr{L} + \frac{\partial \mathscr{L}}{\partial \Psi_\ell} \left( \delta \Psi_\ell + \partial_\nu \Psi_\ell \delta x^\nu \right) + \mathscr{O}(\delta^2)
    \end{align*}
\end{itemize}
带回(\refeq{3.2}):
\begin{equation*}
  \begin{aligned}
\Delta I &= \int d^4 x \left\{ 
\Delta \mathscr{L} + \partial_\mu F^\mu 
+ \partial_\mu \delta x^\mu 
\mathscr{L}[\Psi_\ell (x'), \partial_\mu \Psi_\ell (x')] 
\right\} \\
&= \int d^4 x \left[\frac{\partial\mathscr{L}}{\partial \Psi_\ell}-\frac{\partial \mathscr{L}}{\partial (\partial_\mu \Psi_\ell)} \right]\delta\Psi_\ell+\partial_\mu\left(\frac{\partial \mathscr{L}}{\partial (\partial_\mu \Psi_\ell)}\delta\Psi_\ell\right)\\
&\quad + \partial_\mu F^\mu 
+ \partial_\mu \delta x^\mu \left[ \mathscr{L} + \frac{\partial \mathscr{L}}{\partial \Psi_\ell} \left( \delta \Psi_\ell 
+ \partial_\nu \Psi_\ell \delta x^\nu \right) \right] 
+ \mathscr{O}(\delta^2)
\Big\} \\
&= \int d^4 x \left\{ 
\left[ \frac{\partial \mathscr{L}}{\partial \Psi_\ell} 
- \partial_\mu \frac{\partial \mathscr{L}}{\partial (\partial_\mu \Psi_\ell)} \right] \delta \Psi_\ell
+ \partial_\mu \left[ F^\mu +\left(\frac{\partial \mathscr{L}}{\partial (\partial_\mu \Psi_\ell)}\delta\Psi_\ell\right)
+ \delta x^\mu \mathscr{L} \right] 
\right\}
\end{aligned}
\end{equation*}
如果在要求作用量在该变换下严格保持不变，同时假设运动方程成立的情况下，我们得到：
\begin{equation*}
    \partial_\mu \mathbb{J}^\mu=0\qquad
    \mathbb{J}^\mu = F^\mu +\frac{\partial \mathscr{L}}{\partial (\partial_\mu \Psi_\ell)}\delta\Psi_\ell
+ \delta x^\mu \mathscr{L} 
\end{equation*}
写为更加简洁的形式，\text{\color{blue}诺特流}是：
  \begin{equation}\label{3.3}
      \begin{aligned}
        &\partial_\mu \mathcal{J}^\mu=0\\
        &\mathcal{J}^\mu=\underbrace{F^\mu}_\text{规范项}+\frac{\partial \mathscr{L}}{\partial (\partial_\mu \Psi_\ell)}\delta_{total}\Psi_\ell
        -\delta x^\nu\left\{\frac{\partial \mathscr{L}}{\partial (\partial_\mu \Psi_\ell)}\partial_\nu\Psi_\ell-g^\mu_\nu \mathscr{L} \right\}
      \end{aligned}
\end{equation}
\topictitle{时空平移对称性和能动张量}
在平直时空中，假定四维时空具有平移均匀性，我们因此考虑以下变换：
  \begin{itemize}
  \item \textbf{坐标平移}
  \begin{equation*}
    x^\mu \rightarrow x^\mu + \epsilon^\mu, \quad \epsilon^\mu = \text{const vector}
  \end{equation*}
  \item \textbf{场的不变性}
  \begin{equation*}
    \Psi_\ell (x') \rightarrow \Psi_\ell(x)
  \end{equation*}
\end{itemize}
所以我们得到：
\begin{equation*}
\left\{
  \begin{aligned}
    &\delta_{\text{total}} \Psi_\ell = \Psi_\ell (x') - \Psi_\ell(x) = 0 \\
    &\partial^\nu \Psi_\ell(x) \delta x^\nu = \epsilon^\nu \Psi_\ell(x)
  \end{aligned}
\right.  
\end{equation*}

代入诺特流\refeq{3.3}
\begin{align*}
  \mathbb{J}^\mu &= 
  F^\mu + \frac{\partial \mathscr{L}}{\partial (\partial_\mu \Psi_\ell)} \delta_{\text{total}} \Psi_\ell
  - \delta x^\nu \left\{
  \frac{\partial \mathscr{L}}{\partial (\partial_\mu \Psi_\ell)} \partial_\nu \Psi_\ell
  - g^\mu_{\ \nu} \mathscr{L} \right\} \\
  &= F^\mu - \epsilon^\nu \left\{
  \frac{\partial \mathscr{L}}{\partial (\partial_\mu \Psi_\ell)} \partial_\nu \Psi_\ell
  - g^\mu_{\ \nu} \mathscr{L} \right\}
\end{align*}
不考虑规范变换,剩下的诺特流是：
\begin{equation*}
    \mathbb{J}^\mu = -\epsilon^\nu \left\{
    \frac{\partial \mathcal{L}}{\partial (\partial_\mu \Psi_\ell)} \partial_\nu \Psi_\ell
    - g^\mu_{\ \nu} \mathscr{L} \right\}
\end{equation*}


\noindent
  $\epsilon$的选择具有任意性时,这便是一个守恒量,即:\textbf{能量动量张量}
    \begin{equation}
        \begin{aligned}
          &\partial_\mu \mathcal{T}^\mu_{\ \nu} = 0\\
          &\mathcal{T}^\mu_{\ \nu} = \partial_\nu \Psi_\ell \frac{\partial \mathscr{L}}{\partial (\partial_\mu \Psi_\ell)}
          - g^\mu_{\ \nu} \mathscr{L}
        \end{aligned}
  \end{equation}


\topictitle{Lorentz不变性和角动量张量}
In four-dimensional spacetime, we assume not only homogeneity but also isotropy. Hence, we consider the following transformation:
\begin{itemize}
  \item \textbf{Infinitesimal Lorentz transformation of coordinates}
  \[
    x^\mu \rightarrow x^\mu + \delta\omega^{\mu\nu} x^\nu, \quad \omega_{\mu\nu} = -\omega_{\nu\mu} = \text{constant antisymmetric tensor}
  \]
  \item \textbf{Intrinsic transformation of the field}
  \[
    \Psi_\ell (x') = \Psi_\ell(x) + \delta_{\text{total}} \Psi_\ell(x), \quad
    \delta_{\text{total}} \Psi_\ell = \frac{i}{2} \omega_{\mu\nu} (S^{\mu\nu})^{a}_{\ b} \varphi ^{(b)}
  \]
\end{itemize}
Substituting into the Noether current expression:
\begin{align*}
   \mathbb{J}^\mu &= F^\mu +
   \frac{\partial \mathscr{L}}{\partial (\partial_\mu \Psi_\ell)}
   \delta_{\text{total}} \Psi_\ell
  - \delta x^\nu \left\{ \frac{\partial \mathscr{L}}{\partial (\partial_\mu \Psi_\ell)} \partial_\nu \Psi_\ell
  - g^\mu_{\ \nu} \mathscr{L} \right\} \\
  &= F^\mu +
  \frac{i}{2} \omega_{\nu\rho} (S^{\nu\rho})^{a}_{\ b} \varphi ^{(b)}
  \frac{\partial \mathscr{L}}{\partial (\partial_\mu \Psi_\ell)} +
  \omega^{\nu\rho} x_\rho \mathcal{T}^\mu_{\ \nu}
\end{align*}
We rewrite $\omega^{\nu\rho} x_\rho \mathcal{T}^{\mu\nu}$ by exchanging indices:
\begin{align*}
\omega^{\nu\rho} x_\rho \mathcal{T}^\mu_{\ \nu} 
= \omega^\nu_{\ \rho} x^\rho \mathcal{T}^{\mu\nu}
= \frac{1}{2} \omega_{\nu\rho} \left( x^\rho \mathcal{T}^{\mu\nu} 
- x^\nu \mathcal{T}^{\mu\rho} \right)
\end{align*}
Thus the Noether current becomes:
\begin{align*}
  \mathbb{J}^\mu &= F^\mu +
  \frac{i}{2} \omega_{\nu\rho} (S^{\nu\rho})^{a}_{\ b} \varphi ^{(b)}
  \frac{\partial \mathscr{L}}{\partial (\partial_\mu \Psi_\ell)} +
  \frac{1}{2} \omega_{\nu\rho} \left( x^\rho \mathcal{T}^{\mu\nu} 
  - x^\nu \mathcal{T}^{\mu\rho} \right) \\
  &= F^\mu +
  \frac{1}{2} \omega_{\nu\rho} \left[ i (S^{\nu\rho})^{a}_{\ b} \varphi ^{(b)}
  \frac{\partial \mathscr{L}}{\partial (\partial_\mu \Psi_\ell)} +
   \left( x^\rho \mathcal{T}^{\mu\nu} 
  - x^\nu \mathcal{T}^{\mu\rho} \right) \right]
\end{align*}

\noindent
If we consider only Lorentz invariance and neglect internal symmetries, we obtain:
    \begin{equation}
    \mathbb{J}^\mu =
    \frac{\omega_{\nu\rho}}{2} \left[ i (S^{\nu\rho})^{a}_{\ b} \varphi ^{(b)}
    \frac{\partial \mathscr{L}}{\partial (\partial_\mu \Psi_\ell)} +
    \left( x^\rho \mathcal{T}^{\mu\nu} 
    - x^\nu \mathcal{T}^{\mu\rho} \right) \right]
  \end{equation}
Clearly, this is a conserved quantity. We define:
    \begin{equation}
    \mathcal{M}^{\mu\nu\rho} = i (S^{\nu\rho})^{a}_{\ b} \varphi ^{(b)}
    \frac{\partial \mathscr{L}}{\partial (\partial_\mu \Psi_\ell)} +
    \left( x^\rho \mathcal{T}^{\mu\nu}
    - x^\nu \mathcal{T}^{\mu\rho} \right) \quad \text{and} \quad
    \partial_\mu \mathcal{M}^{\mu\nu\rho} = 0
  \end{equation}
\noindent
It can be seen that this conserved quantity consists of two terms, which are, in fact, the “orbital angular momentum density” and the “spin angular momentum density”:
    \begin{equation}
    \begin{aligned}
          \mathcal{M}^{\mu\nu\rho}&
    = \underbrace{i (S^{\nu\rho})^{a}_{\ b} \varphi ^{(b)} 
    \frac{\partial \mathscr{L}}{\partial (\partial_\mu \Psi_\ell)}}_{\text{Spin angular momentum density } } 
    + \underbrace{\left( x^\rho \mathcal{T}^{\mu\nu} 
    - x^\nu \mathcal{T}^{\mu\rho} \right)}_{\text{Orbital angular momentum density } }\\
    &=\mathcal{M}^{\mu\nu\rho}+\mathscr{L}^{\mu\nu\rho}
    \end{aligned}
  \end{equation}

